\documentclass{notaaula}

        \titulonota{Fenômenos ondulatórios}
        \creditos{Turma 2A · Física · Novembro/2026 · Prof. Josué Cavalcante}

        \begin{document}
        \cabecalho

        \begin{sobre}
        Reflexão, refração, difração, polarização e interferência, com aplicações tecnológicas. Habilidades em foco: EM13CNT106, EM13CNT107 e EM13CNT310.
        \end{sobre}

        \section{Reflexão e refração}

\begin{copiar}{Mudanças na propagação}
Na reflexão, a onda retorna ao meio de origem. Na refração, atravessa a interface e muda de rapidez;
a frequência permanece igual à da fonte. Como $v=\lambda f$, o comprimento de onda muda com $v$.
Em uma corda, a reflexão em extremidade fixa ocorre com inversão; em extremidade livre, sem inversão.
\end{copiar}

\section{Difração e polarização}

\begin{copiar}{Contornar e selecionar}
Difração é o espalhamento da onda ao atravessar uma abertura ou contornar um obstáculo. Ela é mais
perceptível quando o tamanho da abertura é comparável a $\lambda$.

Polarização seleciona uma direção de oscilação e só ocorre com ondas transversais. O fato de a luz
poder ser polarizada é evidência de seu caráter transversal.
\diaadia{Óculos polarizados reduzem reflexos intensos de superfícies horizontais.}
\end{copiar}

\section{Superposição e interferência}

\begin{copiar}{Somar deslocamentos}
Quando ondas ocupam a mesma região, os deslocamentos instantâneos se somam. Depois da superposição,
cada onda continua sua propagação.
\begin{itens}
\item Interferência construtiva: crista com crista, amplitude resultante maior.
\item Interferência destrutiva: crista com vale, cancelamento parcial ou total.
\end{itens}
Para duas fontes coerentes em fase,
\[
  \begin{aligned}
  \Delta r&=m\lambda &&\text{(construtiva)},\\
  \Delta r&=\left(m+\frac12\right)\lambda &&\text{(destrutiva)}.
  \end{aligned}
\]
\end{copiar}

\begin{exemplo}
Duas ondas chegam em fase com amplitudes $3\un{cm}$ e $2\un{cm}$. Na superposição construtiva,
a amplitude instantânea máxima é $\resultado{5\un{cm}}$. Se chegarem em oposição de fase, a
amplitude resultante é $1\un{cm}$.
\end{exemplo}

\section{Ondas estacionárias e ressonância}

\begin{copiar}{Nós e ventres}
Ondas estacionárias surgem da interferência entre ondas de mesma frequência que se propagam em
sentidos opostos. Nós permanecem com amplitude nula; ventres apresentam amplitude máxima.
Em uma corda fixa nas duas extremidades,
\[
  L=n\frac{\lambda}{2},
  \qquad
  f_n=n\frac{v}{2L},\quad n=1,2,3,\ldots
\]
Ressonância ocorre quando uma força periódica tem frequência próxima de uma frequência natural,
produzindo grande amplitude.
\end{copiar}

\begin{exemplo}[Modo fundamental]
Uma corda de $L=\dec{1,2}\un{m}$ tem onda com $v=120\un{m/s}$. No modo fundamental,
$\lambda_1=2L=\dec{2,4}\un{m}$ e
\[
  f_1=\frac{v}{\lambda_1}=\resultado{50\un{Hz}}.
\]
\end{exemplo}

\begin{dica}
Interferência não destrói energia de forma permanente. Regiões de mínimo são compensadas por
regiões de máximo, conforme a distribuição da onda.
\end{dica}

\begin{exercicios}
\nivel{1}{Conceitos}
\begin{questoes}
\item O que permanece constante na refração de uma onda?
\item Quando a difração se torna mais perceptível?
\item Por que ondas longitudinais não podem ser polarizadas?
\item Diferencie interferência construtiva e destrutiva.
\item Defina nó e ventre em uma onda estacionária.
\nivel{2}{Aplicação}
\item Duas ondas em fase têm amplitudes $4\un{cm}$ e $3\un{cm}$. Ache a amplitude resultante máxima.
\item Para $\lambda=\dec{0,80}\un{m}$, uma diferença de caminho $\Delta r=\dec{1,6}\un{m}$ gera qual tipo de interferência?
\item Uma corda fixa de $2\un{m}$ sustenta o modo fundamental. Ache $\lambda_1$.
\item Se $v=200\un{m/s}$ na corda do item anterior, ache $f_1$.
\nivel{3}{Síntese}
\item Explique o funcionamento básico de um fone com cancelamento ativo de ruído.
\item Uma abertura de $\dec{0,50}\un{m}$ difrata mais uma onda de $\lambda=\dec{0,40}\un{m}$ ou de $\lambda=\dec{0,01}\un{m}$? Justifique.
\item Cite uma situação em que a ressonância é útil e outra em que pode ser perigosa.
\end{questoes}
\gabarito{1) frequência; 2) abertura ou obstáculo comparável a $\lambda$;
6) $7\un{cm}$; 7) construtiva, pois $\Delta r=2\lambda$; 8) $4\un{m}$;
9) $50\un{Hz}$; 11) a onda de $\dec{0,40}\un{m}$.}
\end{exercicios}


\end{document}
